\section{Introduction}

\bo{text below should be changed because it has several positioning issues}

\bo{The revised message should be: When an agent successfully reproduces a research result, to what extent does that success imply that it faithfully evaluated the supplied artifact? What kinds of interventions break that correspondence? Can we detect those changes in a systematic way?}

\bo{Or, more precisely: What interventions should count as legitimate artifact operation, and how do we observe when an agent crosses that boundary?}

\bo{For \aebench, justify every decision. One thing to keep in mind is answering these questions ``We need feature/decision X because ...''}

\bo{The issue with current agentic artifact evaluation (and agentic reproducibility, more broadly) is that these frameworks largely check the final result. However, I don't know of any work that makes the point that agent interventions that could help/make the artifact run (e.g., fix missing dependencies, broken environment variables, etc.) may alter measurement conditions. Probably adding missing dependencies is okay and won't affect much, but fixing hardcoded values or even worse -- small bugs in scripts, could.}

\bo{The high-level research question: Can an agent distinguish between tasks where it should intervene to simply satisfy a goal versus tasks where it should preserve what is being observed so that the result remains meaningful?}

Artifact Evaluation (AE) is a community-driven process intended to improve the long-term value of papers accepted at computer systems conferences. Authors submit code, data, and experimental results for review so that others can reproduce, compare, and extend the work. Many venues recognize successful artifact evaluation through typically one or more of three badges: Artifact Available, Artifact Functional, and Results Reproduced. Artifact Available means the artifact is publicly archived. Artifact Functional means reviewers were able to install or exercise it successfully. Results Reproduced means the paper's main results were obtained by someone other than the authors using the provided artifact.
~\bo{This paragraph became too long. I guess my original idea was to say this: 
(1) AE has become an important part of the peer review process; 
(2) it isn't just about checking final results; 
(3) it it actually about committees independently obtaining evidence corresponding to the paper claims (so the process and steps matter as much as the final result)
So, the second half can be deleted}

A core premise of artifact evaluation is to ensure artifacts provide evidence to support the main paper's claims. In practice, however, reviewers often reach that scientific question only after substantial operational work: reconstructing environments, repairing build assumptions, preparing workloads, and producing interpretable outputs.
\bo{This introduces the ``operational burden'' of artifact evaluation, but we should also indicate the issue: committees sometimes need to repair artifacts which actually may not preserve what should be evaluated and might actually perform the work of authors}

Our experience on artifact evaluation committees across several systems conferences is that a significant fraction of the reviewer-author discussion focuses on how to operate the artifact, fix bugs, and provide missing context. For example, in one recent conference with 50 registered artifacts, the evaluation process generated more than 1,600 reviewer-author messages. Of these, 52\% concerned functionality or bug-fixing issues, while only about 15\% concerned whether the artifact and results supported the paper's main claims. Common problems included missing dependencies, hardcoded paths that failed on a new machine, undocumented hardware or software requirements, brittle setup logic, and script bugs that prevented the workflow from completing. These issues are often ancillary to the paper's main contribution, but they determine whether evaluation can or cannot proceed.
\bo{This need more conferences. It could actually become a mini-study since we have concrete data for about 11-12 AE sessions. One interesting information si tu quantify these bug-fixing comments (e.g., config bugs, logic bugs, missing scripts, etc.). I guess we might want to reframe things here from ``AE is hard'' (so what?) to ``what does real artifact evaluation boils down to in practice vs in theory''}

A natural response is to require stricter artifact packaging, much as conferences require strict paper formatting. Authors could be asked to provide exhaustive READMEs, containers, and standardized execution scripts. This can reduce some setup and operational failures, yet not fully solve the problem. Paper formatting standardizes a static document, whereas artifact packaging standardizes only part of an executable system. Docker can make an artifact easier to launch, but it does not guarantee correct functionality, independent reproducibility, or agreement with the paper's claims. Key details often remain outside the container, including dataset preparation, workload availability, dependency versions, hardware assumptions and topology, sequencing requirements, result difference tolerances, and paper-to-code mapping. Moreover, this approach shifts substantial labor onto the authors. Academic incentives reward new publications, not long-term artifact maintenance, so stricter packaging requirements cannot eliminate the need to handle artifacts that remain incomplete, brittle, or under-specified.
\bo{Too long, cut in at least half. And I guess there's another issue here: the main takeaway shouldn't be ``containers don't solve everything'' (like, yeah, obviously), but rather ``packaging can allow for sort-of stalbe/predictable execution, but it doesn't guarantee evaluation-preserving execution'' because successful evaluation still depends on executing the intended experiment under the appropriate conditions.}

Thus, the AE process is largely a sequence of concrete checks. Reviewers must determine whether an artifact can be obtained, built, configured, and executed, and whether its outputs remain consistent with the paper. This staged structure makes artifact evaluation a plausible target for both partial automation and measurement. Recent AI systems are increasingly capable of carrying out multi-step software tasks, including setup, debugging, and execution. This creates an opportunity to treat artifact evaluation as a measurable systems task: agents can be evaluated on the procedural work that reviewers already perform against explicit success conditions.
\bo{I agree here this is now ``old news''. It is the case that agents can now operate and run research prototypes (several papers in this space). However, this should stress that AE differs from ordinary goal orieented software tasks because AE committees must achieve progress without arbitrarily changing the system under evaluation.}

Recent work on agentic artifact evaluation (e.g., Artisan [1]~\bo{need to collect all relevant work, not just Artisan}) shows that AI agents can automate important parts of the reproducibility process. Artisan [1] formulates artifact evaluation as a code generation task. Given a paper, a software artifact, and a table, Artisan generates a reproduction script that can be executed independently to reproduce the target table. Artisan also introduces an automatic judging mechanism that checks whether the generated script produces the expected result and whether this is the result of meaningful evaluation, not simply copying already checked-in results. This is an important step towards making agentic artifact evaluation more auditable and demonstrates that agents can help with specific results reproduction tasks that can be checked by independent validation steps (e.g., oracles).
\bo{I think the opportunity here is to list some of the ``neighboring'' RQs and cite literature besides Artisan. Existing work primarily evaluates success through reconstructed outputs or artifact execution. But AEBench is evaluating exeuction procedure and modification evaluation explicit by checking if agents really follow instructions.}

Artifact evaluation, however, involves a broader set of steps/workflow, and failures in this workflow are often localized: an artifact may be available but not compilable, or compilable but not executable, or executable but unable to reproduce results within a reasonable tolerance. A single end-to-end success label hides these distinctions. For instance, an agent that installs dependencies and builds the artifact but fails to prepare workloads receives the same final failure as one that never successfully establishes the environment. For artifact evaluation assistance, this distinction is important because partial progress can still reduce reviewer burden. We therefore need a benchmark that exposes the procedural stages of artifact evaluation and provides explicit validation for each stage.
\bo{We might get push-back here. Breaking down a large process into smaller stages is not really a novel concept. But it can still be useful/interesting if this paragraph would explain is how this helps tell something interesting about the agents}

Our work starts from a simple premise: agents should be evaluated on the work reviewers actually do. We therefore propose AEBench, a benchmark for artifact evaluation built from artifacts that accompany peer-reviewed CS papers. Each benchmark entry packages the paper, the artifact, the original instructions, and machine-checkable evaluation logic. The goal of AEBench is to measure how far an agent progresses through the artifact evaluation workflow, identify the operational stage that blocks further progress, and record the evidence produced before that point.
\bo{This paragraph comes slightly too late. Also, it is not clear what we're actually saying. I think this paragraph should clearly answer this (1) what does an agent successful on AEBench looks like?; and (2) what does that tell us about the skills/capabilites on said agent (or the agent that fails)?}

Each AEBench entry is built from a research artifact accompanying a paper accepted for publication and already evaluated through an official artifact evaluation process. Specifically, the artifact was awarded all three artifact evaluation badges: Artifact Available, Artifact Functional, and Results Reproduced. The benchmark preserves the original codebase and operational guidelines along with the challenges reviewers face in operating research artifacts, including incomplete instructions, outdated dependencies, changing execution environments, implicit compilation assumptions, brittle datasets or benchmark preparation steps, and result checking that requires clear explanation of expected variance.
\bo{This paragraph is not being clear on what AEBench consists. It sounds like we're preserving artifacts with incomplete instructions, broken dependencies, etc. But it sounds weird becuse these are artifacts that earned all three badges so they (in theory) have none of that.}

AEBench organizes artifact evaluation into four stages: environment setup, artifact build, benchmark preparation, and experiment execution. These stages reflect the dependency structure of artifact evaluation: an agent cannot run experiments before it has built the system, and it cannot build the system before it has established the required environment. Stage-specific evaluation distinguishes an agent that can establish the environment from one that can also build the artifact, prepare datasets or benchmarks, and execute the experiments.
\bo{Probably needs to be bumped one paragraph up. After a few readings of this it starts to sound like we're making claims that  artifacts should always follow these four stages in order. Also ``experiment execution'' contains result validation as well, but that needs to be explicitly included in the description because readers would wonder where is the most important part of artifact evaluation. Plus, these stages are kind of milestones we observed over many artifacts, instead of pre-defined stages set in stone for AE work.}

[workflow figure placeholder]

Figure 1. AEBench artifact evaluation workflow.

AEBench executes each artifact through a single orchestration layer. This layer prepares an isolated workspace, instantiates the target artifact, starts the intended execution environment, and instructs the agent to perform the evaluation task using the artifact's documentation (e.g., README file). It records the agent's actions and the evidence produced during the run. After the agent exits, the same layer runs stage-specific oracles that programmatically check whether the agent satisfied the minimal success conditions for each stage and produce a structured report along with a final score from 0 to 4.
\bo{This paragraph is not being clear that we don't merely assign a numeric score. Two agents can score 4/4 on the same artifact: one actually reproduces the artifact, while the other decides to heavily modify and alter the experiment}

AEBench also includes mechanisms to detect ``shortcut behaviors'' that could satisfy an oracle without actually exercising the artifact. This is particularly important during experiment execution, where an agent could, for example, write an expected result directly or modify a script to emit values already available in the paper. AEBench addresses this in two complementary ways. First, our framework implements a lightweight proxy around the underlying system shell, records the commands executed by the agent together with their outputs, exit status, runtime, and selected filesystem changes, providing evidence about how result files were produced and whether experiment-driving scripts or binaries were modified. Second, after the agent finishes, AEBench discards its generated outputs and replays the submitted setup and execution scripts in a newly provisioned environment controlled by the framework. Only results regenerated during this replay are considered evidence by the evaluation oracles when checking. These mechanisms do not 100\% prove that every reported value was produced by the intended computation, but they make common forms of ``evaluation gaming'', like directly writing expected outputs or modifying the code that generates them, substantially harder to exploit.
\bo{``shorcut behavior'' is vague, but ``cheating'' is too strong; maybe ``gaming''? We should also make explicit that tracing aims to determine what changes did the agent make, and which of them matter to the validity of the evaluation? Also, the reply part is not obsolete. This was some old idea that never got implemented}

AEBench is intended for AI agents meant to offset the procedural burden of artifact evaluation. These agents are intended to support, not replace, the scientific assessment performed by human reviewers. Questions such as whether a reproduced measurement is sufficient to support a claim, whether a deviation is acceptable, or whether the implementation faithfully captures the intended method still require human interpretation. Thus, AEBench evaluates agents on the part of artifact evaluation that is repetitive, repeatable, and machine-checkable, so that reviewer effort can be redirected, with agent support, toward assessing the paper's claims against the evidence produced by the artifact.
